\section{Simulation study}\label{sec:simulation}

\todonote{The tables and figures of this section are currently based on $R=50$ replications per case (the run on our side); the full design with 500 replications is generated by the enclosed code (\texttt{run\_sim.R}, about 7--8 hours on 8 cores), after which \texttt{summarize\_sim.R} and \texttt{update\_manuscript\_sim.py} regenerate all tables and figures of this section and of the Appendix. The Monte Carlo standard errors in the Appendix refer to the actual number of replications; the coverage proportions in Table~\ref{tab:simcoverage} in particular need the full run. The numbers quoted in the text should be re-checked against the regenerated tables.}

As briefly mentioned in Section 2.4, usage of marginal models with incomplete data (FEM2 for frequency and SEM2/SEM3 for severity) could be innocuous in case the deductible can be assumed to be exogenous so that these marginal models with incomplete claim data can be equivalent to the marginal models with complete claim data (FEM1 and SEM1, respectively). In other words, to the extent that such exogeneity assumption of the deductible is violated, reliability of the marginal models with incomplete data can be susceptible. \rev{While it is to be expected that a misspecified (exogenous) model fitted to less informative (censored or truncated) data performs worse than a correctly specified model fitted to complete data, the simulation study serves three purposes that cannot be settled by this qualitative argument alone. First, it verifies that the parameters of the proposed joint models, including the dependence parameters, are recoverable from censored and truncated data under the sequential estimation of Section 2.3, which is not obvious a priori because the truncated frequency likelihood depends on the severity model through $p(d;\bm{x})$. Second, it quantifies the magnitude of the biases incurred by the marginal models under a moderate degree of endogeneity, and shows that they are of practical relevance for ratemaking. Third, it shows that the efficiency loss of the joint models fitted to incomplete data relative to those fitted to complete data is modest, and it disentangles the direct effect of the endogeneity of the deductible on the frequency model from the indirect effect transmitted through $p(d;\bm{x})$ from a biased severity model (via the FEM2G variant).} \rev{Since the deductible of a property insurance policy is in practice chosen from a small menu of amounts, as in the LGPIF data of Section~\ref{sec:empirical}, the deductible is generated as a discrete (ordinal) variable, so that the joint models under study are the discrete-deductible models SDM1--3 and FDM1--2 of Section~\ref{section:exoendo}.} To concretely assess \rev{these points}, we conducted a simulation study where we generated $\{D_i, N_{i},Y_{i,j} |i=1,\ldots, 5000, \, j=1,\ldots, N_i \}$ with the following marginal distributions:
{\color{blue}
$$
\begin{aligned}
& \bm{x}_i = (1, X^{[1]}_{i}, X^{[2]}_{i})' \in \mathbb{R}^{3}, X^{[1]}_{i} \sim \mathcal{N}(0, 1), \ X^{[2]}_{i} \sim \mathcal{BN}(0.6), \\
& D_i \in \{d_1,\ldots,d_5\}=\{500,\, 1000,\, 2500,\, 5000,\, 10000\}, \quad
  \Pr(D_i \le d_k \mid \bm{x}_i) = \frac{\exp(\kappa_k + \tilde{\bm{x}}_i'\bm{\beta}_D)}{1+\exp(\kappa_k + \tilde{\bm{x}}_i'\bm{\beta}_D)}, \quad k=1,\ldots,4, \\
&N_{i} \sim \mathcal{NB}(r,\,\exp(\bm{x}_i'\bm{\beta}_N)/ r), \\
&Y_{i,j} \sim \mathcal{PR}(\exp(\bm{x}_i'\bm{\beta}_Y), \, \,\alpha)
\end{aligned}
$$
}
where \rev{$\mathcal{BN}(0.6)$ denotes a Bernoulli distribution with success probability $0.6$,} \rev{the deductible follows the proportional-odds (cumulative logit) model used for the LGPIF data in Section~\ref{sec:empirical} with thresholds $\kappa_1<\kappa_2<\kappa_3<\kappa_4$ and covariate vector $\tilde{\bm{x}}_i=(X^{[1]}_{i}, X^{[2]}_{i})'$ (the intercept is absorbed in the thresholds; a negative coefficient means that the covariate increases the deductible),} $\mathbb{E}[N_i]=\exp(\bm{x}_i'\bm{\beta}_N)$, \rev{$\mathbb{V}ar[N_i]= \exp(\bm{x}_i'\bm{\beta}_N) \left(1+\exp(\bm{x}_i'\bm{\beta}_N)/r\right)$}, and
{\color{blue}
$$
\begin{aligned}
&\bm{\kappa} = (0.1, 1.0, 1.8, 2.6), \quad \bm{\beta}_D = (-0.5, -0.5), \\
&\bm{\beta}_N = (-0.5, 0.5, 0.2),\\
&\bm{\beta}_Y = (8.5, 0.2, 0.5), \\
& r=1, \, \alpha= 3.
\end{aligned}
$$
}
Note that $\mathcal{PR}$ is Pareto distribution, a special case of the GB2 distribution where $\sigma=\alpha_1=1$ and $\alpha_2=\alpha$. \rev{The thresholds are chosen such that the five deductible levels have unconditional proportions of about 45\%, 20\%, 15\%, 10\% and 10\%, similar to the LGPIF data, and the intercept of $\bm{\beta}_Y$ is chosen such that a substantial fraction of the ground-up losses, about 42\% when the deductible is exogenous, falls below the deductible.} We also define $N^*_i ={N_i}$ under censored reporting, $N^*_i =\sum_{j=1}^{N_i} I(Y_{i,j}>D_i)$ under truncated reporting\rev{, and 500 datasets are generated for each of the following cases}:

\begin{itemize}
    \item Case 1 \\
    : $D$ is entirely exogenous, \rev{$D_i \perp N_i$, and $D_i \perp Y_{i,j}$}.
    \item Case 2 \\
    : $D\perp Y$, but $(D_i, N_i)$ are paired with a Gaussian copula with $\rho_N<0$.
    \item Case 3 \\
    : $D\perp N$, but $(D_i, Y_{i,j})$ are paired with a Gaussian copula with $\rho_Y<0$.
    \item Case 4 \\
    : $(D_i, N_{i}), (D_i, Y_{i,j})$ are paired with Gaussian copulas with $\rho_N,\rho_Y<0$.
    \item \rev{Case 5 \\
    : as Case 4, but the conditional independence assumption $N\perp Y\,|\,D$ of Section~\ref{sec:condind} is violated: a single latent risk propensity drives the deductible, the frequency and the severity (see below).}
    \item \rev{Case 6 \\
    : as Case 5, with an additional direct dependence between the frequency and the severity given the deductible.}
\end{itemize}
\rev{In Cases 1--4 the data are generated exactly from the model of Section 2.2 with a discrete deductible. Writing $(F_D(d_{k-1};\bm{x}_i),F_D(d_k;\bm{x}_i)]$ for the ``atom'' of the level $d_k$ (with $F_D(d_0;\bm{x}_i)=0$), the deductible level $D_i$ is first drawn from the ordinal model; then a uniform position $U^{N}_i$ is drawn in the atom of $D_i$ and $N_i=F_{N}^{-1}(V^{N}_i;\bm{x}_i)$ with $(U^{N}_i,V^{N}_i)$ distributed according to the Gaussian copula $C_{DN}$; likewise, for each claim an independent position $U^{Y}_{i,j}$ is drawn in the same atom and $Y_{i,j}=F_{Y}^{-1}(V^{Y}_{i,j};\bm{x}_i)$ with $(U^{Y}_{i,j},V^{Y}_{i,j})\sim C_{DY}$. Integrating the uniform positions over the atom gives precisely $f_{N|D}$ and $f_{Y|D}$ of \eqref{eq:FDM}--\eqref{eq:SDM}, and $N_i$ and $Y_{i,j}$ are conditionally independent given $(D_i,\bm{x}_i)$. In Case 5 the three positions $U^{N}_i$ and $U^{Y}_{i,j}$ are replaced by the single latent uniform $U^{D}_i$ that generates the level $D_i$ itself, i.e., $D_i=F_D^{-1}(U^{D}_i;\bm{x}_i)$: the pairs $(D_i,N_i)$ and $(D_i,Y_{i,j})$ still follow $C_{DN}$ and $C_{DY}$, and $N_i\perp Y_{i,j}\,|\,U^{D}_i$, but $N_i$ and $Y_{i,j}$ are dependent given the observed level $D_i$ through their common position inside the atom (policy-years with many claims tend to sit at the end of the atom with the smaller latent deductible and larger losses). In Case 6, in addition, the latent normal variable of each $Y_{i,j}$ has correlation $\rho_{NY}=0.3$ with the residual of the latent normal variable of $N_i$ given the deductible, which induces a positive dependence between the frequency and the severity given the deductible. Cases 5 and 6 thus assess the sensitivity of the joint models to two forms of violation of the conditional independence assumption.} \rev{The dependence parameters are $\rho_N=-0.5$ in Case 2, $\rho_Y=-0.5$ in Case 3, and $(\rho_N,\rho_Y)=(-0.6,-0.4)$ in Cases 4--6; they are the parameters of the Gaussian copulas between the latent uniforms and the claims, and because of the ties in the ordinal deductible the Kendall's $\tau$ between the observed level and the claims is attenuated (for example, in Case 4 the average Kendall's $\tau$ between $D_i$ and $N_i$ is about $-0.28$ and that between $D_i$ and $Y_{i,j}$ about $-0.15$, compared with $2\arcsin(\rho)/\pi=-0.41$ and $-0.26$ on the latent scale).}

Through the sequential maximum likelihood scheme as in the paragraphs following  \eqref{eq:jointloglik}, we estimated the parameters \rev{$\bm{\kappa}$, $\bm{\beta}_D$ of the ordinal deductible model, $\bm{\beta}_Y$, $\alpha$ (and $\rho_Y$) of the severity models SEM1--3 and SDM1--3, and $\bm{\beta}_N$, $r$ (and $\rho_N$) of the frequency models FEM1, FEM2, FEM2G and FDM1--2} under each case with the reporting scenarios\rev{; the joint models use the Gaussian copula, i.e., the copula family of the data-generating process. As in Section~\ref{sec:empirical}, FEM2 and FEM2G evaluate $p(d;\bm{x})$ with the severity models SEM3 and SEM1, respectively, and FDM2 with the joint severity model SDM3 and the fitted deductible model. Starting values are obtained from the data of the respective scenario (a log-linear regression for the severity, a Poisson regression for the frequency and $\rho=0$ for the copula); the true values are never used.} The estimation performance for parameter $\theta$ was assessed with the following measures:
$$
\begin{aligned}
\text{Mean}&:= \overline{\hat{\theta}}=\frac{1}{500} \sum_{r=1}^{500} \hat{\theta}^{(r)}, \\
\text{RMSE}&:= \sqrt{\frac{1}{500} \sum_{r=1}^{500} (\hat{\theta}^{(r)} - \theta)^2}, \\
\text{Bias}&:= \overline{\hat{\theta}} - \theta, \\
\text{Var}&:= \frac{1}{500} \sum_{r=1}^{500} (\hat{\theta}^{(r)} - \overline{\hat{\theta}} )^2,
\end{aligned}
$$
where $\hat{\theta}^{(r)}$ is the estimated value of $\theta$ under the specified model using the $r${th} simulation sample. \rev{In the tables of the Appendix, the entry ``(s.e.)'' below each mean denotes the Monte Carlo standard error of the mean, $\sqrt{\text{Var}/500}$, so that a bias that is small relative to this Monte Carlo standard error is compatible with an unbiased estimator; note also that $\text{RMSE}^2=\text{Var}+\text{Bias}^2$. Since these are finite-sample Monte Carlo results for one particular design, we refer to an estimator whose bias is negligible relative to its RMSE as ``approximately unbiased'' in what follows, without claiming exact unbiasedness.} \\
To assess the predictive performance of insurance claims with $\hat{\theta}^{(r)}$, the following mean absolute errors (MAE) of \rev{the frequency and severity components of the pure premium in} (\ref{eq:censoredpremium}), (\ref{eq:truncatedpremium}) are measured under each model,
\begin{align*}
    \text{Frequency}&:=\frac{1}{500}\sum_{r=1}^{500}\left(\frac{1}{5000}\sum_{s=1}^{5000}|N^{*(r)}_s -\hat{N}^{*(r)}_s|\right),\nonumber\\
    \text{Severity}&:=
        \frac{1}{500}\sum_{r=1}^{500}\left(\frac{1}{S^{(r)}}\sum_{\rev{t}=1}^{S^{(r)}}|Y^{*(r)}_{\rev{t}} -\hat{Y}^{*(r)}_{\rev{t}}|\right),\nonumber
\end{align*}
 where $S^{(r)}$ is the total number of \rev{claim} observations under censored reporting and truncated reporting in $r$th dataset and $\hat{N}^{*(r)}_s$ / $\hat{Y}^{*(r)}_{\rev{t}}$ are predicted values of claim frequency / severity in $r$th dataset, respectively\rev{, namely the conditional means of the observed frequency and severity given the deductible and the covariates under the fitted model: $\hat{\rm E}(N|D=d;\bm{x})$ and $\hat{\rm E}(\{Y-d\}_+|D=d;\bm{x})$ under censored reporting, and $\hat{\rm E}(N^*|D=d;\bm{x})=\hat p(d;\bm{x})\hat{\rm E}(N|D=d;\bm{x})$ and $\hat{\rm E}(Y-d|Y>d,D=d;\bm{x})$ under truncated reporting, where the marginal models evaluate ${\rm E}(N|D=d;\bm{x})$ and the distribution of $Y$ given $D=d$ without conditioning on the deductible. Because these MAEs are dominated by the irreducible variability of the claims (in particular of the heavy-tailed severities), we also report the error of the premium components themselves, which is what matters for ratemaking:}
{\color{blue}
\begin{align*}
    \text{Frequency}_{\text{prem}}&:=\frac{1}{500}\sum_{r=1}^{500}\left(\frac{1}{5000}\sum_{s=1}^{5000}|m^{(r)}_s -\hat{m}^{(r)}_s|\right),\qquad
    \text{Severity}_{\text{prem}}:=\frac{1}{500}\sum_{r=1}^{500}\left(\frac{1}{S^{(r)}}\sum_{t=1}^{S^{(r)}}|\mu^{(r)}_t -\hat{\mu}^{(r)}_t|\right),
\end{align*}
}
\rev{where $m_s^{(r)}$ and $\mu_t^{(r)}$ are the true conditional means of the observed frequency and severity given $(D,\bm{x})$ under the data-generating process (the frequency and severity factors of \eqref{eq:censoredpremium} or \eqref{eq:truncatedpremium}, computed with the true parameters) and $\hat m_s^{(r)}$, $\hat\mu_t^{(r)}$ their estimates under the fitted model. The detailed estimation results (mean, Monte Carlo standard error, RMSE and bias of each parameter) are reported in Tables~\ref{tab:AppSimDed}--\ref{tab:AppSimEndCase6} in the Appendix, and are summarized graphically in Figures~\ref{fig:simcase1}--\ref{fig:simcase4}; with one exception (the deductible model in one replication of Case 6, whose line search terminated abnormally at estimates that are nevertheless regular), the optimizer converged in all replications and for all models.}

@@CASE_DISCUSSION@@

\FloatBarrier

